%% Compile with XeLaTeX (Overleaf: Menu -> Compiler -> XeLaTeX).
\documentclass[conference]{IEEEtran}
\usepackage{iftex}
\RequireXeTeX   % Overleaf: Menu -> Compiler -> XeLaTeX (pdfLaTeX cannot typeset the Hindi quotations)
\usepackage{fontspec}
\setmainfont{texgyretermes}[Extension=.otf, UprightFont=*-regular, BoldFont=*-bold, ItalicFont=*-italic, BoldItalicFont=*-bolditalic]
% Devanagari font for the Hindi: the bundled file if present (zip upload), else a font Overleaf already has.
\IfFileExists{fonts/NotoSerifDevanagari-Regular.ttf}{%
  \newfontfamily\devfont[Path=fonts/, Script=Devanagari, Scale=0.92]{NotoSerifDevanagari-Regular.ttf}}{%
  \IfFontExistsTF{Noto Serif Devanagari}{%
    \newfontfamily\devfont[Script=Devanagari, Scale=0.92]{Noto Serif Devanagari}}{%
    \newfontfamily\devfont[Script=Devanagari]{FreeSerif}}}
\newcommand{\dev}[1]{{\devfont #1}}
\usepackage{amsmath,amssymb}
\usepackage{booktabs,tabularx,array}
\usepackage[hyphens]{url}
\usepackage[hidelinks]{hyperref}
\newcolumntype{Y}{>{\raggedright\arraybackslash}X}
\renewcommand{\arraystretch}{1.12}
\setlength{\tabcolsep}{3.5pt}
\begin{document}

\title{Assessing Legal Meaning Preservation in English and Hindi Simplifications of Indian Statutes}
\author{\IEEEauthorblockN{Karan Singh, Atharv Arora, Jatin Khanchi\\ S24CSEU1939, S24CSEU2439, S24CSEU1932}\IEEEauthorblockA{School of Computer Science Engineering and Technology\\ Bennett University, Greater Noida, India\\ \textit{CSET 346 Natural Language Processing (project PID-08), instructor Dr. Surabhi Solanki}}}
\maketitle
\begin{abstract}
Simplifying legal text increases accessibility. However, if simplifications do not maintain legal meaning (i.e. removing conditions, changing legal obligations to permissions, or altering time frames), they may be misleading. JUDGEBERT introduced a learned metric for legal meaning preservation (LMP) on French insurance text. Here, we create such a metric for legal text in Indian legislation written in English and Hindi. Using text from The Right to Information Act, 2005 and the Consumer Protection Act, 2019, we generate 338 simplified and original legal text pairs. In one of the styles there is a specified type of legal error introduced by simplification. These are annotated according to the JUDGEBERT metric (1-10) by two disclosed large language models. As such, all annotations here are performed by large language models, as opposed to humans, a deviation we report throughout the paper. Training a cross encoder on this data alone gives poor results (Pearson = 0.228 \ensuremath{\pm} 0.288) and is not reliably better than the best general purpose baseline (calibrated LaBSE 0.101). After pre-training on FrJUDGE, our main XLM-R model achieves Pearson=0.576 \ensuremath{\pm} 0.102. A post-hoc MuRIL variant matches it (0.577 \ensuremath{\pm} 0.068 averaged over 5 seeds). This model generalises better in Hindi and between different acts. However, when we re-run pre-training on FrJUDGE, we see dependence on the pre-training run. Zero-shot English-Hindi transfer remains close to 0. A variant-mean lookup that uses the answer key reaches 0.831. For this reason, we report correlations within variants. Finally, for the optional legal error type, training a MuRIL classifier with data including injected errors achieves macro-F1 of 0.789 compared to 0.509 for rules. Injected negations are handled far better than by the embedding baselines, while obligation-to-permission flips are handled poorly. No system passes both sanity checks.
\end{abstract}
\begin{IEEEkeywords}
legal NLP, text simplification, meaning preservation, evaluation metrics, Hindi, multilingual transformers, cross-encoder regression
\end{IEEEkeywords}

\section{Introduction}

Statutory text in India, and in both English and Hindi, is long and complex. Large language models can rewrite such text into plain language, which is attractive for access to justice at scale. The risk is that the rewriting is judged by the wrong yardstick. General similarity measures reward lexical and semantic overlap, so they give high scores to a simplification that deletes a proviso but keeps nearly every other word, even though the legal consequence has changed.

In [1], it was shown that simple general similarity measures fail to capture legal meaning. On FrJUDGE, 297 French insurance sentence blocks scored by five law students, BERTScore and SBERT achieved Pearson correlations with human LMP scores of 0.46 and 0.13 respectively. However, a CamemBERT-v2-based cross encoder trained directly on the scoring achieved 0.74. BERTScore and MeaningBERT over-score the human label on 82\% of test pairs, while the learned measure never does. This demonstrates that LMP should be specifically learned, rather than approximated with general similarity measures.

We aim to determine whether these results generalise to another legal system, a different pair of languages, and a different kind of statutory text. Additionally, we address the shortcomings of the previous study. For each shortcoming, we make contributions:

\begin{itemize}
\item JUDGEBERT only covers French and a single domain. We include legal text from 2 acts in central Indian legislation, from different domains. Within one act (Right to Information Act), we have both English and Hindi text (as the official Hindi text in the other act could not be found (see Section III-A)). Where available in both languages, the Hindi text is a direct translation of the same provision in English. Therefore, comparisons between languages are controlled.
\item Only one simplification style exists in FrJUDGE. We introduce aggressive simplifications, which often remove conditions. We also include perturbations with an additional legal error of known type. This means we have a diagnostic dataset to understand what types of legal error a metric cannot detect.
\item As inter-annotator agreement on the LMP label in FrJUDGE was low (Krippendorff’s Alpha = 0.10), with two groups: lenient and strict, agreement is considered a significant result. We include agreement from the actual population of annotators used.
\item FrJUDGE splits are made at pair level, so a test original may reappear in training as an identical augmentation pair, and the high correlation (0.97 Pearson) for the augmented model was calculated on a test set consisting of many trivial identical and unrelated pairs. In contrast, our splits are grouped by original clause, and have been confirmed to be disjoint at clause and text level. We include separate sanity checks on hold-out provisions that are used nowhere else.
\item JUDGEBERT reports no out-of-domain test. We train on one act and test on the other.
\end{itemize}

Our contributions include: i) Bilingual LMP dataset with a documented, blinded labelling protocol, and explicitly described label producers. ii) Multilingual LMP cross encoder following the JUDGEBERT approach. iii) Comparison with general metrics. In place of the LLM judge that our labelling decision ruled out, we make a transfer learning setup by pre-training on FrJUDGE and fine-tuning on our data. iv) Cross linguistic, cross domain and sanity checks. v) Classifier for detecting legal error type, with evaluation on precision, recall and F1 score. vi) Qualitative error analysis. In Section VII we include two examples of failure within our pipeline, which would cause confidently incorrect results. Both would be difficult to detect without specific testing.

\section{Related Work}

\subsection{Evaluating simplification}

Simplification quality is typically assessed in terms of simplicity, fluency, and meaning. Common approaches to measuring meaning include n-gram similarities, contextual embeddings (BERTScore [3]), sentence-level cosine similarity, or learned or question-based measures (QuestEval [9], LENS [10], MeaningBERT [11]). BLESS [2] evaluates large language models as sentence simplifiers using zero shot prompts. We use the same prompt to ensure our normal variant is comparable to existing practices.

\subsection{Legal meaning preservation}

This work follows from JUDGEBERT [1]. The scoring criteria (which we follow identically) scores simplifications as Accurate (7-10), Imprecise (2-6), or Off track (1) and subtracts 1 for each identified error (with a minimum score of 1). Error types are omission, hallucination, inconsistency, and confusion. The model is a regression cross encoder on concatenated text. Augmented data includes identical pairs (score 10) and unrelated pairs (score 0). Unrelated data is taken from Quebec legal texts with ROUGE score < 0.25 and BLEU score < 25.

\subsection{Indian legal NLP}

Existing work on Indian legal text focuses on court judgements. Domain adapted models (InLegalBERT [12]) and benchmarks (IL-TUR [13]) allow tasks such as judgement prediction, rhetorical role labelling, and summarisation. Statutory simplification, and whether a simplification is legally faithful, is comparatively unexplored in Indian law. We are unaware of any Hindi resource for it.

\subsection{Multilingual encoders}

We consider several models. XLM-R [5] is multilingual, but not focused on Indic languages. MuRIL [6] and IndicBERT v2 [7] are trained to handle Indian languages and transliteration. LaBSE [4] provides high quality language agnostic sentence embeddings and is included as a baseline.

\section{Dataset}

\subsection{Sources and provenance}

We select two central acts for this corpus. They are from different legal domains, and have been selected as official translations exist in Hindi. These are the Right to Information Act, 2005 (public information) and the Consumer Protection Act, 2019 (consumer law). India Code is the canonical repository for both, but during construction of this corpus version we were unable to access it. Instead, we sourced text from other official government resources. We note exactly where we sourced from. See Table I.

\begin{table}[!t]
\centering
\caption{Sources, editions and extraction yield}
\label{tab:1}
\footnotesize
\begin{tabularx}{\columnwidth}{YY>{\centering\arraybackslash}p{33pt}>{\centering\arraybackslash}p{21pt}}
\toprule
\textbf{Document} & \textbf{Edition} & \textbf{Sections} & \textbf{Units} \\
\midrule
RTI Act 2005, English (DoPT) & As on 18 Nov 2025; includes the 2023 amendment & 31 & 109 \\
RTI Act 2005, Hindi (DoPT) & Carries the 2019 amendments; sec. 8(1) is the older wording; text layer unusable, read by OCR (Tesseract 5.4, tessdata\_best hin) & 27 & 101 \\
Consumer Protection Act 2019, English (Gazette) & Act 35 of 2019, 9 Aug 2019; original enactment & 107 & 229 \\
Consumer Protection Act 2019, Hindi & Not used: unavailable when the corpus was built (Bill text as introduced found later) & -- & -- \\
\bottomrule
\end{tabularx}
\end{table}

As a result of this, there are two implications in this work. Firstly, we were unable to find the official Hindi version of the Consumer Protection Act. Therefore, this act in this corpus only contains English text. The English-Hindi comparison (E2) is restricted to the Right to Information Act. Secondly, the English and Hindi versions of the Right to Information Act do not correspond to the same edition. The English version (current to 18 November 2025) includes 2023 amendments. However, in the Hindi version, whilst the 2019 amendments appear in footnotes, section 8(1) contains outdated wording for personal information. Provisions affected by this are flagged automatically. During review, such provisions were omitted. To ensure we do not produce fabricated statutory text by modifying one language to match the other, we do not silently reconcile discrepancies.

\subsection{Extraction}

Each PDF is split into sub-section units, such as section 7(1), with the marginal or inline heading included. English and Hindi versions of the same provision get the same key, which makes the bilingual alignment automatic. In real government documents three issues arose that are worth recording, since each one cost us a parser bug. Sub-section references (e.g. (3) of Section 8) should not be parsed as new sub-sections. Text from amendment notes, footnotes, and the structure table (at the beginning of the act) should be removed. However, ‘Repeal’ should not be confused with the abbreviation for footnotes, ‘Rep.’. In the Gazette version of the Consumer Protection Act, section headings are in the margin. Therefore, to extract this document we have implemented a geometry-based backend to separate the marginal section headings and citations from the main text. If any numbered section is found without its heading, the backend throws an error.

\subsection{The Hindi text layer}

Although the Hindi document is not scanned, its text layer cannot be used as it is, and the way it fails is instructive. When extracted using Poppler (pdftotext), 4.9\% of Devanagari characters were lost and replaced with the replacement character (U+FFFD). In total, 2461 Devanagari characters across 100 sub-sections (out of 102) were lost, and this loss is irreversible. When extracted using PyMuPDF, no characters were lost. However, the legacy characters used in the document were mapped to incorrect Unicode blocks (Vedic and Sundanese scripts instead of Devanagari). This resulted in 2681 characters being represented by only 48 unique code points. Additionally, vowels were ordered visually rather than logically. Whilst readable by humans, this text was incorrectly structured for machines. To avoid extracting such text, we calculate a percentage of Devanagari characters in the text, and a percentage of incorrectly ordered vowels. If the text fails either of these checks, it is not accepted. For both methods detailed, this occurs, so both are correctly rejected. This is important to note as if text was silently accepted in either case, then a corpus of seemingly readable but nonsensical Hindi would result. Since the text appears readable, no downstream metrics would indicate any issue with it.

To extract the text from the Hindi document, we perform optical character recognition on the document. The page was scaled to 300dpi, then extracted using Tesseract v5.4 with the best quality hindi model (tessdata\_best). This model was preferred over the fast model in this study (based on previous literature and our own trials) due to improved performance on legal terminology. This means that the text passes our criteria extremely well (a Devanagari ratio of 0.993, a misplaced-vowel ratio of 0.000 and zero replacement characters). However there is still a systematic issue with the OCR in this document, and it matters because it hits exactly the part of a statute that must not be approximated. Specifically, the number 1 is often missed or incorrectly identified. As a result section 12 becomes 2, section 13 becomes 03, and (1) sub-sections become empty brackets. To correct this, we repair the numbering from the sequence rather than from how it looks. If the number in the text is a known 1-digit omission of the expected number based on the sequence, we attempt to correct it. Otherwise, we leave it unchanged. All corrections are recorded.

\subsection{Clause selection}

Provisions were eligible if they contained between 1-5 sentences, 15-90 English words, and a Flesch readability score of <=50 to ensure a non-trivial simplification task. Additionally, they were ineligible if they were boilerplate text (short titles, extent, definitions, rule making powers, repeals and savings), and if they were repeated in either act. It is very important to remove repeated provisions. Repeated provisions with identical text are common in legislation, and having two identical provisions across different splits (despite being grouped by clause) would lead to leakage from test data to train data. A maximum of three clauses from each section were chosen. Provisions containing exceptions, overrides, conditions and time limits were weighted in favour of selection. This is where the legal meaning is most delicate. To account for loss during review, provisions were oversampled 1.3 times.

Every candidate then passes a review step, which is the only place where a clause can be accepted. The selection tool never makes the keep decision itself; it only pre-computes flags identifying potential issues (replacement characters, Latin script in Hindi text, numeric discrepancy, missing Hindi text, known edition difference). In the design, the review step is performed by a human. However, due to time constraints, this was completed using a rule-based agent script. In this script, rows with misalignment, incomplete content, Latin script in Hindi text or within edition difference were discarded, with reasons recorded for each. The only correction allowed was number corrections in the way described above. Only corrections where the number in the Hindi text differs by 1 (dropped by OCR) from the English text were made. Seven corrections were made. They were recorded, along with the original and corrected text. A sample of kept rows was then checked against the source pages. No agent at any point edited the statutory wording.

The corpus contains 78 clauses. Of these, 38 are from the Right to Information Act, 35 of which have corresponding Hindi text. 40 clauses are from the Consumer Protection Act, and the Hindi text was unavailable. Therefore, 113 clause-language items are present in the corpus. There are 3 variations per item, minus one perturbation for which the rule found nothing to edit. There are therefore 338 pairs in the corpus. The brief requires 300 pairs with human judgements. As this corpus contains LLM labels, this requirement is not satisfied. (See Section III-D.) Table II shows breakdown of this data.

\begin{table}[!t]
\centering
\caption{Corpus composition by act, language and variant}
\label{tab:2}
\footnotesize
\begin{tabularx}{\columnwidth}{YY>{\centering\arraybackslash}p{29pt}>{\centering\arraybackslash}p{21pt}>{\centering\arraybackslash}p{21pt}>{\centering\arraybackslash}p{21pt}}
\toprule
\textbf{Act} & \textbf{Lang} & \textbf{Clauses} & \textbf{Pairs} & \textbf{Mean words} & \textbf{Mean label} \\
\midrule
CPA & EN & 40 & 119 & 50.2 & 7.40 \\
RTI & EN & 38 & 114 & 53.0 & 7.69 \\
RTI & HI & 35 & 105 & 52.5 & 7.49 \\
All &  & 78 & 338 & 51.9 & 7.53 \\
\bottomrule
\end{tabularx}
\end{table}

\subsection{Simplifications}

Each selected clause is given 3 variations in each language. The normal variation is given the zero-shot BLESS prompt, translated for Hindi. The aggressive variation asks for an explanation for a child aged 12, with a maximum of 2 sentences. This reliably sacrifices conditions, and so provides examples from the lower end of the scale not provided by FrJUDGE. The perturbed variation is a simplified text, which has had one rule-based legal-edit applied, balanced across five types: modal changes (shall/must \ensuremath{\rightarrow} may, and vice versa), negation, quantity change, removal of exceptions/ provisos/ overrides, and restriction of scope. For the Hindi text, edits can only be made to the last verb phrase. For example, karega becomes kar sakega. Unrelated words that merely contain the same letters are not altered. The idiom ‘as may be prescribed’ has been explicitly excluded from editing, since the may in this phrase confers no permission. The base simplification is not guaranteed to be error free. Therefore, perturbed pairs have 1 extra error introduced. However, the number of errors compared to the original legal text may not be exactly 1. For example, the normal simplified text for section 12(6) of the RTI Act in Hindi already narrows the term vidhan-mandal (legislature) to vidhansabha (assembly). The perturbed text keeps this error as well.

Although this pipeline limits text lengths to 200 tokens (English) and 500 tokens (Hindi), as opposed to 100 tokens (JUDGEBERT), it was not used for this corpus. One difference to note is that JUDGEBERT’s simplifications were sampled from the gpt-4-turbo model, using temperature = 1.0 and top-p = 0.9. Simplifications in this corpus were generated by a Claude Opus 5.5 model, following the same prompts. There were no sampling parameters to record. However, the prompt ID, model and generation time are recorded. Temperature and top-p fields contain null values. Therefore, the simplifications generated here are likely to be more conservative than if sampled from temperature=1.0, which plausibly shifts the label distribution upward compared to FrJUDGE. Since no sampling API was involved, no output cap was enforced. In the imported data, the max\_tokens (null) and finish\_reason (always written as stop) fields are not observations of how generation ended.

\subsection{Scoring protocol and its deviation from human judgement}

In this version of the work the protocol has deviated, and we state it plainly. Instead of 3 human annotators (with common calibration), two disclosed LLM annotators were used due to time constraints. Both annotations were performed by Claude Opus 5.5 models, and each completed a blinded workbook according to the same rubric and conventions as the human annotators. This therefore does not satisfy the requirement of agreement with human legal judgment. None of the claims in this paper should be assumed to do so.

Blinding was implemented mechanically for what each annotator could open, and it is only partial. Each annotator is given only the anonymised pair ID. For each annotator, pairs are in random order. Pairs for the same clause are separated. The variant name, error type and answer are in separate files which were never viewed by the annotator. Neither annotator saw the other annotator’s workbook. The scoring was not directly input. Instead, the bracket was selected and errors were counted. From this, scores were calculated using our own code. Therefore, a rubric change cannot be lost in a spreadsheet formula. However, since models were used as annotators, several issues arise. Firstly, since models from the same families were used in generation and annotation, annotators cannot be considered blinded if they annotated text they generated. Since this is the case in this experiment, it is disclosed. Secondly, evaluating an LLM judge against labels from its own model is not allowed, as it measures self-consistency rather than preservation of meaning. Since both annotators in this experiment are from the same model, this rule removed the LLM-judge baseline entirely. Any judge we could access would have been scored against labels from its own model family. No score from an LLM judge is therefore reported. We record its absence as a consequence of the labelling choice, rather than quietly dropping a planned comparison. This is the first experiment to be run once human labels are available.

There are 3 notes to consider here. Firstly, since both annotators are from the same model (Claude Opus 5.5), agreement between two models from the same family was calculated, rather than two independent judges. Agreement is likely inflated by shared pre-training. Secondly, the level of blinding differed between annotators. Annotator 1 was given only the anonymised IDs and the text. They knew general information about corpus design, but did not know which items contained injected errors. Annotator 2 was used to generate the simplifications. Although the variant label in their workbook was hidden, they may have recognised their own output, and thus deduced which versions had perturbations. Therefore, this annotator was in fact not blinded. Thirdly, since there were 2 annotators, the leave-one-out statistic is simply the Pearson correlation between them. However, this is not an upper bound on prediction of mean scores. Each annotator has a Pearson correlation of 0.904 and 0.949 to the mean label (including their own score). Therefore, this statistic is only useful to show agreement between the LLMs, and no conclusions can be made about what proportion of maximum agreement this represents.

Krippendorff’s \ensuremath{\alpha} [8] was 0.610 for the interval LMP scores and 0.503 for the nominal brackets. Inter-LLM agreement had Pearson correlation 0.723. Of the 338 pairs, 68 appear in the adjudication list, and none of them has been adjudicated. These figures are far above the 0.10 agreement that JUDGEBERT reports between its 5 law students, but the comparison is lopsided: two instances of one model agreeing is a much weaker result than 5 humans agreeing. So the gap tells us about model consistency, not about the construct being well defined. Furthermore, the annotators are systematically offset. The average scores for the two models were very different (8.115 and 6.935). Thus, in some circumstances, the mean label may be a numerical compromise between two very different labelling schemes rather than an adjudicated judgement. The distribution of labels is bimodal, concentrated around 10 and a tail from 1 to 5. Individual labels rarely occur in 6 or 7 because the rubric subtracts from bracket maxima of 10, 6 and 1, but mean labels do (15 pairs at 6 and 7 at 7). This bimodal distribution affects RMSE. One convention differed between the two annotators, and we state it rather than smoothing it over. If inserting/removing a negation caused a reversal in the whole operation of the provision, one model assigned Off-track, while the other assigned Imprecise with confusion. The mean label with the two values (as with FrJUDGE) was assigned. Every such divergence between the two annotators is recorded on the adjudication list.

To keep a path to the original claim open, a separate validation set was constructed in a 30 pair workbook for human legal annotators. It was stratified by language and variant. This workbook has yet to be annotated. Therefore, there is no human validation of these labels. As such, the correlation between our model labels and human legal judgement is unmeasured. Consequently, all results in this paper should be considered with this in mind. The workbook and aggregation tools can be found in the repository. Measuring this value simply requires annotation of this workbook.

\subsection{Augmentation, sanity holdouts and splits}

As with JUDGEBERT, for each clause and language pair we add one identical pair (labelled 10) and one unrelated pair (where the paired provision has ROUGE-1, ROUGE-2 and ROUGE-L scores <0.25 and BLEU scores <25). Unrelated pairs have label 1, rather than 0, since 0 cannot be achieved and using 0 would bias the regression. A flag reproduces the original choice. Overlap statistics were calculated using a Devanagari aware tokenizer. Using the commonly-used rouge\_score package would cause all Hindi pairs to be classified as unrelated (since non-ASCII text is ignored), and so would bias the augmented set.

Sanity holdouts were built only from provisions that are not selected, are not in a section containing a selected clause, and are not used as augmentation partners, so passing these checks cannot be an artefact of memorising training data. 60/10/30 splits were made by clause across 3 seeds (42-44). All pairs generated from the same clause, both languages, all three variants and augmented partners will all be in the same split. This is guaranteed by the splitter at the clause and text level. Headline metrics only consider annotated test pairs. Augmented pairs and sanity check pairs are excluded and reported separately.

\section{Method}

\subsection{LMP model}

As with JUDGEBERT, pairs were embedded by concatenating original, delimiter and simplified text. This embedding was produced using a pre-trained multilingual model. A linear regression head on the pooled representation predicts the score and is trained with MSE loss. Predictions were clipped to the 1-10 rubric. XLM-R and MuRIL models (IndicBERT v2 was intended but not executed) were trained with and without augmentation, with 3 seeds (42-44) grouped by clause. Checkpoints were selected using the annotated development set. Our main system (FrJUDGE-pretrained XLM-R fine-tuned with augmentation) was fixed in advance. For experiments E1-E4, the same configuration was used. All systems used are listed in tables in the same order. The featured system in each table was not selected based on performance on test data.

Some changes were made to the published parameters due to a property of our data, and we state both. In JUDGEBERT, learning rate = 5e-5, batch size = 16, and training stopped early with patience = 5. We have learning rate = 2e-5, batch size = 8, no early stopping, and 30 training epochs (saving best dev checkpoint). Linear learning rate decay and gradient clipping (norm 1.0) remain unchanged. Only \textasciitilde{}35 annotated pairs exist in our dev set, so the dev loss is dominated by noise. When using early stopping with patience=5 (or 10), the model reached early stopping and predicted the mean label. This resulted in near 0 correlation. The same issue occurs in the original data in Section VII during reproduction of FrJUDGE (patience=5 caused 1 in 10 seeds to fail). This issue is due to the early stopping interacting badly with a small dev set, and not just our data. We disclose that these parameters were chosen after observing this issue in training. This is a form of selection on the dev set, and so it is recorded as a limitation.

\subsection{Baselines}

Four reference points were compared, in increasing order of sophistication. ROUGE-L and BLEU, using our Devanagari aware tokenizer, capture direct overlap. LaBSE [4] captures language agnostic meaning of sentences with cosine similarity. Multilingual BERTScore [3] F1 captures contextual token matches, and is the most effective general purpose metric in the original study. Finally, a fourth baseline was designed but was intentionally not run. An LLM judge was to be prompted with the exact rubric given to annotators (temperature 0). This would investigate if a general model given the rubric in its prompt performs similarly to a model trained on it. However, as explained in Section III-F, this baseline is inapplicable with our data. Since all labels in this study come from the same model family, this baseline would measure model self-consistency. The code for this baseline has been created, but was not run. For all other baselines, 2 scores were calculated. First, score was calculated naively, and then with a linear calibration fitted on the training split only. Both are shown, as uncalibrated scores may give good correlation but meaningless RMSE.

\subsection{Decisions and reporting}

For the decision output, the continuous predictions were mapped onto the rubric brackets. Scores of 7+ were Preserved, scores of 2-7 were Partially Preserved and scores of less than 2 were Not Preserved. In addition to Pearson and RMSE, we report Spearman, MAE and over-rating (proportion of pairs predicted higher than actual label). Since in practice, it is dangerously misleading for a metric to be high for a simplification missing a condition, we include this statistic. All overall metrics are means, with standard deviations given for 3 seeds. Results are sub-grouped by language, act and variant.

\section{Experiments}

Five experiments are defined in advance. E1: Comparison between generic baselines and trained cross-encoders on random clause-grouped splits. The LLM judge belonged here and is reported as not run (see Section III-F). E2: Per-language results and zero-shot transfer in both directions. Training on English, testing on Hindi, and vice versa, is possible due to equivalence in provisions. Only Right to Information Act evaluated (see Section III-A). E3: Cross-act training and testing both ways. E4: Sanity checks on the disjoint holdouts. Proportion of identical pairs with scores \ensuremath{\geq}9.9, and unrelated pairs with scores \ensuremath{\leq}1.1. E5: Qualitative analysis of at least 15 difficult examples, categorised by the type of introduced error and the most important change as noted by each annotator. The perturbation allows for systematic analysis.

All steps from annotated pairs to the following tables form a reproducible pipeline. Files are created at each stage (splitting, baselines, training, evaluation and error analysis), and all tables in this paper are derived from the resulting files.

\begin{table*}[!t]
\centering
\caption{E1: baselines and LMP model on annotated test pairs}
\label{tab:3}
\footnotesize
\begin{tabularx}{\textwidth}{Y>{\centering\arraybackslash}p{58pt}>{\centering\arraybackslash}p{58pt}>{\centering\arraybackslash}p{54pt}>{\centering\arraybackslash}p{54pt}>{\centering\arraybackslash}p{46pt}}
\toprule
\textbf{System} & \textbf{Pearson} & \textbf{Spearman} & \textbf{RMSE} & \textbf{MAE} & \textbf{Over-score \%} \\
\midrule
XLM-R, FrJUDGE\ensuremath{\rightarrow}ours, DA & 0.576 \ensuremath{\pm} 0.102 & 0.544 \ensuremath{\pm} 0.106 & 2.037 \ensuremath{\pm} 0.142 & 1.474 \ensuremath{\pm} 0.213 & 38.8 \ensuremath{\pm} 3.9 \\
XLM-R, FrJUDGE\ensuremath{\rightarrow}ours, no DA & 0.586 \ensuremath{\pm} 0.098 & 0.553 \ensuremath{\pm} 0.101 & 2.037 \ensuremath{\pm} 0.110 & 1.525 \ensuremath{\pm} 0.042 & 46.1 \ensuremath{\pm} 7.0 \\
XLM-R from scratch, DA & 0.228 \ensuremath{\pm} 0.288 & 0.259 \ensuremath{\pm} 0.280 & 2.464 \ensuremath{\pm} 0.224 & 2.020 \ensuremath{\pm} 0.469 & 47.9 \ensuremath{\pm} 12.1 \\
XLM-R from scratch, no DA & -0.030 \ensuremath{\pm} 0.048 & -0.020 \ensuremath{\pm} 0.044 & 2.444 \ensuremath{\pm} 0.090 & 2.245 \ensuremath{\pm} 0.079 & 45.8 \ensuremath{\pm} 5.6 \\
MuRIL from scratch, DA & 0.067 \ensuremath{\pm} 0.097 & 0.097 \ensuremath{\pm} 0.060 & 2.607 \ensuremath{\pm} 0.180 & 2.365 \ensuremath{\pm} 0.135 & 41.7 \ensuremath{\pm} 6.6 \\
MuRIL from scratch, no DA & -0.042 \ensuremath{\pm} 0.014 & -0.034 \ensuremath{\pm} 0.028 & 4.443 \ensuremath{\pm} 0.058 & 3.876 \ensuremath{\pm} 0.084 & 6.4 \ensuremath{\pm} 2.1 \\
labse-cal & 0.101 \ensuremath{\pm} 0.026 & 0.224 \ensuremath{\pm} 0.066 & 2.417 \ensuremath{\pm} 0.124 & 2.184 \ensuremath{\pm} 0.090 & 44.4 \ensuremath{\pm} 4.8 \\
labse & 0.101 \ensuremath{\pm} 0.026 & 0.224 \ensuremath{\pm} 0.066 & 2.811 \ensuremath{\pm} 0.157 & 2.110 \ensuremath{\pm} 0.121 & 41.0 \ensuremath{\pm} 3.6 \\
bertscore-cal & 0.055 \ensuremath{\pm} 0.040 & 0.115 \ensuremath{\pm} 0.033 & 2.427 \ensuremath{\pm} 0.120 & 2.206 \ensuremath{\pm} 0.088 & 44.7 \ensuremath{\pm} 4.9 \\
bertscore & 0.055 \ensuremath{\pm} 0.040 & 0.115 \ensuremath{\pm} 0.033 & 2.974 \ensuremath{\pm} 0.235 & 2.180 \ensuremath{\pm} 0.232 & 59.2 \ensuremath{\pm} 1.6 \\
bleu-cal & 0.018 \ensuremath{\pm} 0.021 & 0.088 \ensuremath{\pm} 0.042 & 2.433 \ensuremath{\pm} 0.120 & 2.212 \ensuremath{\pm} 0.087 & 45.4 \ensuremath{\pm} 5.8 \\
bleu & 0.018 \ensuremath{\pm} 0.021 & 0.088 \ensuremath{\pm} 0.042 & 5.512 \ensuremath{\pm} 0.238 & 4.977 \ensuremath{\pm} 0.316 & 6.4 \ensuremath{\pm} 1.3 \\
rougeL-cal & 0.010 \ensuremath{\pm} 0.033 & 0.071 \ensuremath{\pm} 0.040 & 2.430 \ensuremath{\pm} 0.119 & 2.217 \ensuremath{\pm} 0.090 & 45.4 \ensuremath{\pm} 5.8 \\
rougeL & 0.010 \ensuremath{\pm} 0.033 & 0.071 \ensuremath{\pm} 0.040 & 4.041 \ensuremath{\pm} 0.159 & 3.563 \ensuremath{\pm} 0.169 & 21.2 \ensuremath{\pm} 1.8 \\
Variant-mean lookup (uses answer key; diagnostic, not a system) & 0.831 \ensuremath{\pm} 0.063 & n/a & 1.385 \ensuremath{\pm} 0.138 & 0.868 \ensuremath{\pm} 0.045 & 33.7 \ensuremath{\pm} 16.7 \\
\bottomrule
\end{tabularx}
\end{table*}

\begin{table*}[!t]
\centering
\caption{E2: results by language, and zero-shot transfer}
\label{tab:4}
\footnotesize
\begin{tabularx}{\textwidth}{YY>{\centering\arraybackslash}p{58pt}>{\centering\arraybackslash}p{54pt}>{\centering\arraybackslash}p{42pt}}
\toprule
\textbf{System} & \textbf{Setting} & \textbf{Pearson} & \textbf{RMSE} & \textbf{Over-score \%} \\
\midrule
XLM-R, FrJUDGE\ensuremath{\rightarrow}ours, DA & train EN+HI, test EN & 0.165 \ensuremath{\pm} 0.063 & 2.380 \ensuremath{\pm} 0.161 & 36.7 \ensuremath{\pm} 8.8 \\
XLM-R, FrJUDGE\ensuremath{\rightarrow}ours, DA & train EN+HI, test HI & -0.041 \ensuremath{\pm} 0.144 & 2.638 \ensuremath{\pm} 0.334 & 45.6 \ensuremath{\pm} 3.8 \\
XLM-R, FrJUDGE\ensuremath{\rightarrow}ours, DA & train EN \ensuremath{\rightarrow} test HI (zero-shot) & 0.035 \ensuremath{\pm} 0.027 & 2.564 \ensuremath{\pm} 0.293 & 42.2 \ensuremath{\pm} 1.9 \\
XLM-R, FrJUDGE\ensuremath{\rightarrow}ours, DA & train HI \ensuremath{\rightarrow} test EN (zero-shot) & 0.131 \ensuremath{\pm} 0.115 & 2.519 \ensuremath{\pm} 0.282 & 31.1 \ensuremath{\pm} 6.9 \\
XLM-R from scratch, DA & train EN+HI, test EN & -0.059 \ensuremath{\pm} 0.197 & 2.391 \ensuremath{\pm} 0.141 & 38.9 \ensuremath{\pm} 7.7 \\
XLM-R from scratch, DA & train EN+HI, test HI & -0.058 \ensuremath{\pm} 0.108 & 2.314 \ensuremath{\pm} 0.146 & 46.7 \ensuremath{\pm} 3.3 \\
XLM-R from scratch, DA & train EN \ensuremath{\rightarrow} test HI (zero-shot) & -0.059 \ensuremath{\pm} 0.086 & 2.583 \ensuremath{\pm} 0.315 & 44.4 \ensuremath{\pm} 3.8 \\
XLM-R from scratch, DA & train HI \ensuremath{\rightarrow} test EN (zero-shot) & 0.028 \ensuremath{\pm} 0.027 & 2.433 \ensuremath{\pm} 0.074 & 37.8 \ensuremath{\pm} 6.9 \\
\bottomrule
\end{tabularx}
\end{table*}

\begin{table*}[!t]
\centering
\caption{E3: cross-act generalisation}
\label{tab:5}
\footnotesize
\begin{tabularx}{\textwidth}{YY>{\centering\arraybackslash}p{54pt}>{\centering\arraybackslash}p{54pt}>{\centering\arraybackslash}p{54pt}}
\toprule
\textbf{System} & \textbf{Split} & \textbf{Pearson} & \textbf{RMSE} & \textbf{MAE} \\
\midrule
XLM-R, FrJUDGE\ensuremath{\rightarrow}ours, DA & CPA \ensuremath{\rightarrow} RTI & 0.388 \ensuremath{\pm} 0.021 & 2.391 \ensuremath{\pm} 0.101 & 1.703 \ensuremath{\pm} 0.031 \\
XLM-R, FrJUDGE\ensuremath{\rightarrow}ours, DA & RTI \ensuremath{\rightarrow} CPA & 0.119 \ensuremath{\pm} 0.051 & 2.629 \ensuremath{\pm} 0.068 & 2.212 \ensuremath{\pm} 0.041 \\
XLM-R from scratch, DA & CPA \ensuremath{\rightarrow} RTI & 0.186 \ensuremath{\pm} 0.053 & 2.603 \ensuremath{\pm} 0.163 & 2.042 \ensuremath{\pm} 0.162 \\
XLM-R from scratch, DA & RTI \ensuremath{\rightarrow} CPA & 0.056 \ensuremath{\pm} 0.110 & 2.589 \ensuremath{\pm} 0.036 & 2.299 \ensuremath{\pm} 0.069 \\
\bottomrule
\end{tabularx}
\end{table*}

\begin{table*}[!t]
\centering
\caption{E4: sanity checks on disjoint holdouts}
\label{tab:6}
\footnotesize
\begin{tabularx}{\textwidth}{Y>{\centering\arraybackslash}p{46pt}>{\centering\arraybackslash}p{46pt}>{\centering\arraybackslash}p{46pt}Y}
\toprule
\textbf{System} & \textbf{Identical \ensuremath{\geq} 9.9 \%: EN (n=50)} & \textbf{Identical \ensuremath{\geq} 9.9 \%: HI (n=3)} & \textbf{Unrelated \ensuremath{\leq} 1.1 \%: EN (n=17)} & \textbf{Unrelated \ensuremath{\leq} 1.1 \%: HI (n=0)} \\
\midrule
XLM-R, FrJUDGE\ensuremath{\rightarrow}ours, DA & 66.7 \ensuremath{\pm} 57.7 & 66.7 \ensuremath{\pm} 57.7 & 94.1 \ensuremath{\pm} 5.9 & not evaluated \\
XLM-R, FrJUDGE\ensuremath{\rightarrow}ours, no DA & 91.3 \ensuremath{\pm} 5.0 & 66.7 \ensuremath{\pm} 0.0 & 17.6 \ensuremath{\pm} 10.2 & not evaluated \\
XLM-R from scratch, DA & 33.3 \ensuremath{\pm} 57.7 & 33.3 \ensuremath{\pm} 57.7 & 31.4 \ensuremath{\pm} 27.8 & not evaluated \\
labse-cal & 0.0 \ensuremath{\pm} 0.0 & 0.0 \ensuremath{\pm} 0.0 & 0.0 \ensuremath{\pm} 0.0 & not evaluated \\
labse & 100.0 \ensuremath{\pm} 0.0 & 100.0 \ensuremath{\pm} 0.0 & 0.0 \ensuremath{\pm} 0.0 & not evaluated \\
bertscore & 100.0 \ensuremath{\pm} 0.0 & 100.0 \ensuremath{\pm} 0.0 & 0.0 \ensuremath{\pm} 0.0 & not evaluated \\
bleu & 100.0 \ensuremath{\pm} 0.0 & 100.0 \ensuremath{\pm} 0.0 & 11.8 \ensuremath{\pm} 0.0 & not evaluated \\
rougeL & 100.0 \ensuremath{\pm} 0.0 & 100.0 \ensuremath{\pm} 0.0 & 0.0 \ensuremath{\pm} 0.0 & not evaluated \\
\bottomrule
\end{tabularx}
\end{table*}

\begin{table}[!t]
\centering
\caption{Mean absolute error on perturbed pairs by injected error type}
\label{tab:7}
\footnotesize
\begin{tabularx}{\columnwidth}{Y>{\centering\arraybackslash}p{21pt}>{\centering\arraybackslash}p{54pt}>{\centering\arraybackslash}p{54pt}}
\toprule
\textbf{Injected error type} & \textbf{n per test split} & \textbf{MAE XLM-R, FrJUDGE\ensuremath{\rightarrow}ours, DA} & \textbf{MAE labse-cal} \\
\midrule
exception\_drop & 4 \ensuremath{\pm} 2 & 1.673 \ensuremath{\pm} 0.474 & 2.254 \ensuremath{\pm} 0.466 \\
modality & 9 \ensuremath{\pm} 2 & 3.160 \ensuremath{\pm} 0.878 & 2.496 \ensuremath{\pm} 0.334 \\
negation & 8 \ensuremath{\pm} 2 & 1.896 \ensuremath{\pm} 0.533 & 4.629 \ensuremath{\pm} 0.531 \\
number & 5 \ensuremath{\pm} 2 & 2.274 \ensuremath{\pm} 0.466 & 2.722 \ensuremath{\pm} 0.172 \\
scope & 7 \ensuremath{\pm} 2 & 1.860 \ensuremath{\pm} 0.517 & 2.672 \ensuremath{\pm} 0.208 \\
\bottomrule
\end{tabularx}
\end{table}

\begin{table*}[!t]
\centering
\caption{Within-variant behaviour of the principal system on the E1 random splits}
\label{tab:8}
\footnotesize
\begin{tabularx}{\textwidth}{Y>{\centering\arraybackslash}p{25pt}>{\centering\arraybackslash}p{46pt}>{\centering\arraybackslash}p{54pt}>{\centering\arraybackslash}p{58pt}>{\centering\arraybackslash}p{58pt}}
\toprule
\textbf{Variant} & \textbf{n per test split} & \textbf{Mean label} & \textbf{MAE XLM-R, FrJUDGE\ensuremath{\rightarrow}ours, DA} & \textbf{Mean signed error XLM-R, FrJUDGE\ensuremath{\rightarrow}ours, DA} & \textbf{Within-variant Pearson XLM-R, FrJUDGE\ensuremath{\rightarrow}ours, DA} \\
\midrule
normal & 33 \ensuremath{\pm} 1 & 9.96 \ensuremath{\pm} 0.02 & 1.015 \ensuremath{\pm} 0.630 & -0.995 \ensuremath{\pm} 0.640 & 0.314 \ensuremath{\pm} 0.363 \\
aggressive & 33 \ensuremath{\pm} 1 & 7.64 \ensuremath{\pm} 0.27 & 1.205 \ensuremath{\pm} 0.134 & -0.683 \ensuremath{\pm} 0.294 & 0.487 \ensuremath{\pm} 0.063 \\
perturbed & 33 \ensuremath{\pm} 1 & 5.04 \ensuremath{\pm} 0.60 & 2.203 \ensuremath{\pm} 0.057 & 1.961 \ensuremath{\pm} 0.189 & 0.584 \ensuremath{\pm} 0.047 \\
\bottomrule
\end{tabularx}
\end{table*}

We state the from-scratch result before the detail, because on its own it is a negative result and burying it would be worse than reporting it. When trained only on the current corpus, the JUDGEBERT-style cross-encoders do not reliably improve over generic similarity measures. The best performing model (XLM-R with augmentation) has a mean Pearson correlation of 0.228 (s.d. 0.288). This varies widely between seeds (0.56, 0.08, 0.05). With 3 seeds, the model shows a mean correlation of around 1.4 standard errors away from zero, meaning that it is not reliably distinguishable from zero correlation, or from the correlation of LaBSE (0.101). Without augmentation, XLM-R performs worse (-0.030). MuRIL with augmentation performs slightly better (0.067) but without augmentation performs very poorly (RMSE 4.4). This dataset does not demonstrate that this task is impossible to learn, but the model does not reliably learn it from \textasciitilde{}200 training pairs per split. The results of reproducing FrJUDGE (Section VII) and the transfer experiment (below) are consistent with this explanation of insufficient data size, though do not conclusively prove it, as FrJUDGE differs in terms of language, domain, and label distribution.

The per-language breakdown of E1 shows where the weakness lies. Training from scratch gives a correlation of 0.297 for English pairs and 0.010 for Hindi. Pre-training on French gives a main model correlation of 0.675 \ensuremath{\pm} 0.104 for English and 0.308 \ensuremath{\pm} 0.115 for Hindi. This analysis may be misleading as English contains pairs from both acts, whereas Hindi only contains Right to Information pairs. The same runs score 0.669 \ensuremath{\pm} 0.069 on pairs from the Consumer Protection Act and 0.485 \ensuremath{\pm} 0.144 on pairs from the Right to Information Act. Hindi remains the weaker language by a wide margin. Since there are 105 pairs of text recovered using OCR, it is difficult to determine how much of this performance drop is due to OCR noise, sample size, or poor representation of the legal register specific to Hindi legislation in XLM-R.

E2 proper uses only the 35 Right to Information clauses common to both languages, so the languages are compared on the same provisions. Around 126 training pairs were used in each split, from 21 clauses. Neither model learns from this subset. When trained on both languages, the transferred model scores 0.165 \ensuremath{\pm} 0.063 on English pairs, but performs worse on Hindi pairs (-0.041 \ensuremath{\pm} 0.144). Zero-shot transfers in both directions give scores of 0.035 \ensuremath{\pm} 0.027 (English \ensuremath{\rightarrow} Hindi) and 0.131 \ensuremath{\pm} 0.115 (Hindi \ensuremath{\rightarrow} English). Training from scratch gives zero or lower correlations in both directions. E2 is therefore inconclusive about cross-lingual transfer: with a training set this small, a null result cannot distinguish a language gap from a data gap. These numbers are also lower than the per-language E1 figures, where the Hindi pairs are trained alongside all English pairs of both acts.

Cross-act generalisation (E3) is asymmetric. Training on Consumer Protection Act pairs and testing on Right to Information Act pairs gives a correlation of 0.388 \ensuremath{\pm} 0.021 (with transfer) and 0.186 \ensuremath{\pm} 0.053 (no transfer). The reverse direction gives a correlation of 0.119 \ensuremath{\pm} 0.051 (with transfer) and 0.056 \ensuremath{\pm} 0.110 (no transfer). Augmentation partners in these splits come only from the training act. Therefore, no text from the target act is seen in training. The two directions differ in more than the domain. In the first case, models are trained on 40 English clauses and tested on both languages. In the second case, models are trained on 38 clauses from both languages and tested on English. With only one pair of acts, we cannot determine what difference has caused this asymmetry.

All generic baselines perform very poorly. LaBSE gives a correlation of 0.101 \ensuremath{\pm} 0.026, multilingual BERTScore gives 0.055 \ensuremath{\pm} 0.040, BLEU 0.018 \ensuremath{\pm} 0.021, and ROUGE-L 0.010 \ensuremath{\pm} 0.033. This is a large drop compared to the performance on FrJUDGE (BERTScore = 0.46), and we attribute this mainly to the current corpus. As a third of pairs have an artificial legal error, the simplified version is almost semantically and lexicographically identical to an accurate example, but with a different meaning. An overlap or embedding metric is not designed to see that difference. This is clearly visible in the diagnostic table. LaBSE's error is largest exactly where the change is smallest on the surface. Negation inversion examples have a mean absolute error of 5.896 \ensuremath{\pm} 0.372, compared to dropping exceptions, with a mean absolute error of 3.022 \ensuremath{\pm} 0.910. (These figures use the naive LaBSE mapping; the diagnostic table uses the train-calibrated mapping, so its values differ.) Correlations as low as this may also suggest an issue with implementation of the baselines. As we did not repeat the baselines on FrJUDGE, we cannot be certain of this. However, all four metrics fail in the same way. Their ranking tracks surface sensitivity as expected.

A surprising finding was the calibration of these models. Under naive mapping, 100\% of identical examples in the holdout sets have a score of 9.9 or higher in all 4 models. However, none of the models give scores of 1.1 or less in unrelated pairs (BLEU gives 11.8\%). Therefore, all generic models fail the second sanity check. Under linear calibration based on training data, all four models fail both sanity checks, achieving 0\%. Calibration causes the distributions of predicted labels to become centred on the mean label value. This reduces RMSE but causes models to fail to correctly predict both extremes that a meaning preservation model should capture. Therefore, calibrated RMSE values for these models should be treated as scaling artefacts, not as evidence of fitness. For this reason, results are presented for both mappings. The trained models perform better, but none pass both sanity checks. XLM-R with augmentation from scratch gives scores of 9.9 or higher for 33\% of identical pairs, and scores of 1.1 or lower for 31\% of unrelated pairs. Transferred models have opposing results. With no augmentation, 91.3\% \ensuremath{\pm} 5.0 of identical pairs are given scores \ensuremath{\geq}9.9, but only 17.6\% of unrelated pairs. With augmentation however, 94.1\% \ensuremath{\pm} 5.9 of unrelated pairs are given scores \ensuremath{\leq}1.1, but only 66.7\% of identical pairs (s.d. 57.7). Two seeds pass all identical pairs, and one seed passes none. No system we tested passes both checks, as JUDGEBERT-DA did, and the instability of the main system on the identical-pair check is a real defect. Due to almost all holdouts being in English (3 identical pairs in Hindi, 0 unrelated pairs in Hindi), very little information can be drawn regarding performance on Hindi.

The diagnostic set is where the case for a learned metric is clearest, and the pattern is the opposite of what we expected. Negation reversals have a much larger average absolute error than any other errors for the calibrated LaBSE model (4.629 \ensuremath{\pm} 0.531) compared to average absolute error values of approximately 2.3-2.7. Whereas inserting a ‘not’ completely reverses the meaning of the legal sentence, the difference in the sentence embedding is negligible. The main model performs better at negation reversals (1.896 \ensuremath{\pm} 0.533) but has the worst performance at modality errors (3.160 \ensuremath{\pm} 0.878). These results align with the idea that the learned metric can pick up on differences in meaning that similarity metrics cannot. However as these error types are artificially generated, it could be that the model is simply picking up on the change in the template. The weakness on modality is the natural next target, since reversing obligation to permission is a common real failure in simplified legal text.

Overall correlation on this corpus rewards telling the three variants apart. If each pair of test data is predicted based on the mean value of training data from the same variant (which would not be possible in application as this requires knowledge from the answer key) then a correlation of 0.831 \ensuremath{\pm} 0.063 is achieved. This is higher than all learned models. Therefore the within variant correlation is a fairer comparison. Within each variant, the main model has a positive correlation with the labels. In aggressive pairs and perturbed pairs this is 0.487 \ensuremath{\pm} 0.063 and 0.584 \ensuremath{\pm} 0.047 respectively, so it is not merely recognising the variant. However, a low correlation in normal pairs (0.314 \ensuremath{\pm} 0.363) gives little insight as the values are almost constant (9.96). Systematic bias is present. Faithful simplifications receive a score that is around 1 point too low, whereas perturbed pairs receive scores that are around 2 points too high. This is the over-scoring failure, on exactly those pairs where it matters most.

\subsection{Transfer from FrJUDGE}

Because the from-scratch model is unstable, we add an experiment that the original work does not have. First, the cross-encoder is pre-trained on the FrJUDGE dataset (same task and criteria as this dataset, but in French insurance legal text) giving a Pearson correlation of 0.707 with our code. Then it is fine-tuned on the Indian data splits. The scores of these transferred models are provided below, alongside the from-scratch models so the benefit of transfer learning can be seen as an ablation study. Both augmented and un-augmented versions of the transferred model are presented, as well as augmented and un-augmented from-scratch models.

Pre-training with augmentation gives a correlation between 0.85 and 0.87 across our two pre-training runs on the test set of FrJUDGE. Fine-tuning from there with augmentation, our main system gives correlations across the three seeds of 0.564, 0.684 and 0.481 (0.576 \ensuremath{\pm} 0.102). Without augmentation, correlations of 0.625, 0.658 and 0.474 (0.586 \ensuremath{\pm} 0.098) are obtained. As this version with augmentation mirrors that of JUDGEBERT-DA and is used in E2 to E4, it will be labelled as the main model. This was decided after seeing the results of E1, however, the two versions are within error margins of each other. Compared to the from-scratch model (0.228 \ensuremath{\pm} 0.288), the gain is large in level and also in stability. The transferred models have correlations between 0.48 and 0.68 in seeds, whereas the from-scratch models have correlations between 0.05 and 0.56.

This demonstrates that pre-training on the same criteria in a different language gives a large and more stable gain on these labels. However, it does not demonstrate that preserving meaning in legal text is independent of language/jurisdiction. Pre-training in French also means the model learns the scoring range and distribution of labels. The reference labels are judgements from LLMs. In addition, E2 finds no cross-lingual transfer within our own corpus. Further, the transferred model is not yet ready for deployment. The lookup outperforms the transferred model, it fails one sanity check on one seed in three, and in section VI we can see that it is making very poor predictions in certain cases.

\subsection{Follow-up: MuRIL and five seeds}

After the experiments above we ran two follow-ups, both post hoc and reported as such. First, the same transfer was attempted, but using MuRIL [6] instead of XLM-R. This model is pre-trained on Indian languages and was run with 5 seeds (42-46) grouped by clause. Second, a new run of the XLM-R transfer was run with the same 5 seeds. This meant that a new run of pre-training on FrJUDGE had to be completed. Results for these are shown in Table IX. The main model remains as-is. Promoting a better model after seeing these numbers would be selection on the test data.

\begin{table*}[!t]
\centering
\caption{Follow-up: FrJUDGE\ensuremath{\rightarrow}ours with augmentation, Pearson per experiment (mean \ensuremath{\pm} std over seeds; FrJUDGE = test Pearson of the pre-training run)}
\label{tab:9}
\footnotesize
\begin{tabularx}{\textwidth}{Y>{\centering\arraybackslash}p{54pt}>{\centering\arraybackslash}p{54pt}>{\centering\arraybackslash}p{54pt}>{\centering\arraybackslash}p{54pt}>{\centering\arraybackslash}p{58pt}>{\centering\arraybackslash}p{54pt}>{\centering\arraybackslash}p{54pt}}
\toprule
\textbf{System (FrJUDGE\ensuremath{\rightarrow}ours, DA)} & \textbf{FrJUDGE} & \textbf{E1} & \textbf{E1 HI} & \textbf{E2 EN} & \textbf{E2 HI} & \textbf{E3 C\ensuremath{\rightarrow}R} & \textbf{E3 R\ensuremath{\rightarrow}C} \\
\midrule
XLM-R, original pre-training (3 seeds) & 0.873 / 0.849 & 0.576 \ensuremath{\pm} 0.102 & 0.308 \ensuremath{\pm} 0.115 & 0.165 \ensuremath{\pm} 0.063 & -0.041 \ensuremath{\pm} 0.144 & 0.388 \ensuremath{\pm} 0.021 & 0.119 \ensuremath{\pm} 0.051 \\
XLM-R, re-run pre-training (5 seeds) & 0.691 & 0.486 \ensuremath{\pm} 0.112 & 0.290 \ensuremath{\pm} 0.068 & 0.155 \ensuremath{\pm} 0.094 & 0.066 \ensuremath{\pm} 0.141 & 0.339 \ensuremath{\pm} 0.046 & 0.147 \ensuremath{\pm} 0.040 \\
MuRIL (5 seeds) & 0.839 & 0.577 \ensuremath{\pm} 0.068 & 0.466 \ensuremath{\pm} 0.092 & 0.269 \ensuremath{\pm} 0.090 & 0.182 \ensuremath{\pm} 0.119 & 0.394 \ensuremath{\pm} 0.054 & 0.251 \ensuremath{\pm} 0.186 \\
\bottomrule
\end{tabularx}
\end{table*}

The MuRIL transfer model generalises best. It has an overall correlation equal to the main model (0.577 \ensuremath{\pm} 0.068 compared to 0.576 \ensuremath{\pm} 0.102) with less variance. On all generality measures, it performs better. It performs better on Hindi (E1) (0.466 vs 0.308), and in the matched bilingual subset (0.269 and 0.182 vs 0.165 and -0.041). It performs better in cross act pairs (0.251 vs 0.119, though with a standard deviation of 0.186). On zero-shot transfer between languages, every system stays near zero (MuRIL 0.086 (English to Hindi) and 0.104 (Hindi to English)).

By repeating the transfers, we observe a problem with this transfer method. In the XLM-R transfer, with the same seed, pre-training on FrJUDGE gave a lower correlation on the test set of FrJUDGE (0.691 vs 0.873). Consequently, the correlation on the downstream task was also worse. With augmentation, 0.486 \ensuremath{\pm} 0.112, and without augmentation 0.166 \ensuremath{\pm} 0.214. These results are much worse than those shown in Table III for this version of the model (0.586 \ensuremath{\pm} 0.098). Thus performance without augmentation relies on a single favourable run of pre-training. The augmented variants, and MuRIL, are the robust ones.

\subsection{Legal-error type classification}

The brief asks, optionally, for the type of legal error in each pair. We frame it as classification on the standard and perturbed pairs of the E1 data splits, with 6 classes (none, modality, negation, number, omission of exception, scope). The targets are objective: none for standard simplifications, and the injection type for perturbed simplifications, not the LLM annotations. Note that aggressive pairs do not have a known type, and so were omitted. Rule-based baselines are produced using the same cues as given in the cue table in the demonstration. Cross-encoder based classifiers were also created, trained with class weighted cross entropy, keeping the best development macro F1.

\begin{table*}[!t]
\centering
\caption{Legal-error type classification on the E1 splits (five seeds): macro-F1 (mean \ensuremath{\pm} std), accuracy and per-class F1}
\label{tab:10}
\footnotesize
\begin{tabularx}{\textwidth}{Y>{\centering\arraybackslash}p{54pt}>{\centering\arraybackslash}p{21pt}>{\centering\arraybackslash}p{17pt}>{\centering\arraybackslash}p{17pt}>{\centering\arraybackslash}p{17pt}>{\centering\arraybackslash}p{17pt}>{\centering\arraybackslash}p{21pt}}
\toprule
\textbf{System} & \textbf{Macro-F1} & \textbf{Acc.} & \textbf{Mod.} & \textbf{Neg.} & \textbf{Num.} & \textbf{Exc.} & \textbf{Scope} \\
\midrule
Rules (no training) & 0.509 \ensuremath{\pm} 0.074 & 0.552 & 0.45 & 0.63 & 0.58 & 0.33 & 0.46 \\
Majority (“none”) & 0.111 \ensuremath{\pm} 0.001 & 0.502 & 0.00 & 0.00 & 0.00 & 0.00 & 0.00 \\
MuRIL & 0.362 \ensuremath{\pm} 0.174 & 0.327 & 0.25 & 0.64 & 0.57 & 0.07 & 0.47 \\
MuRIL + augmentation & 0.789 \ensuremath{\pm} 0.101 & 0.754 & 0.66 & 0.85 & 0.97 & 0.79 & 0.79 \\
XLM-R + augmentation & 0.581 \ensuremath{\pm} 0.288 & 0.695 & 0.55 & 0.64 & 0.69 & 0.30 & 0.55 \\
MuRIL + augmentation + cues & 0.573 \ensuremath{\pm} 0.072 & 0.598 & 0.47 & 0.74 & 0.73 & 0.38 & 0.49 \\
\bottomrule
\end{tabularx}
\end{table*}

With training on just the data, the rule-based model outperformed the classifiers (macro-F1: MuRIL 0.362 \ensuremath{\pm} 0.174 vs 0.509 \ensuremath{\pm} 0.074). The small number of pairs (around 130 across 6 classes) leads to poor performance. However, when only the training clauses were augmented with each of the applicable error types (both in the simplified text, and original text) (approximately 630 pairs per split), this changes. MuRIL reaches a macro-F1 of 0.789 \ensuremath{\pm} 0.101 (0.852, 0.874, 0.621, 0.779 and 0.817 across seeds) and the model outperformed the rules in all classes. Adding the cues from the rules to the input did not improve performance (0.573). Since the error injection rules used to create the test data are the same as those used in the rule-based model and in the augmentation, this indicates that the model can detect these error types in the text. It gives no information on errors that may be made by simplifiers naturally, as these have no error type annotation.

\section{Error Analysis}

Scoring the corpus blind made one structure plain, and it organises this section. There are two types of error which cause changes in meaning in simplifications. The first type is loss: a proviso, an override, a subject-to clause, the date from which a limitation period runs, or an exception simply disappears. This is inherent in aggressive simplifications, and the resulting text is honest about being shorter. The second type of error causes inversion. This includes adding/removing negation, turning an obligation into a permission (or vice versa), or narrowing a quantifier. For example, “any political party” becomes “certain political party”, and “every order” becomes “some order”. Inversion is particularly risky, as it does not change the length of the text or make it seem less fluent than a truthful simplification, but tells the reader the inverse of the law. Inversion errors also cannot be detected using overlap and embedding measures. As reported in Section V, LaBSE's error is nearly twice as large on negation reversals as on dropped exceptions, which is precisely backwards from how much the surface changed.

Examples have been chosen automatically from annotated pairs and the blinded annotator's own workbook comments. Injected errors were prioritised in order to allow for known error types. Each group is capped so that every error type is shown before any repetition. Within the table for each example, the provision, language, variant, type of change, annotation (mean of annotators), predicted value from each system (with predictions for this pair in test data), and the blinded annotator's own score have been included, so that differences between annotators stay visible rather than hidden inside the mean.

\begin{table*}[!t]
\centering
\caption{E5: legally difficult cases, with the change that was made (pert. = perturbed, aggr. = aggressive, norm. = normal; Annot. = annotator B, the blinded one)}
\label{tab:11}
\footnotesize
\begin{tabularx}{\textwidth}{YYY>{\centering\arraybackslash}p{21pt}>{\centering\arraybackslash}p{25pt}>{\centering\arraybackslash}p{46pt}>{\centering\arraybackslash}p{33pt}>{\centering\arraybackslash}p{37pt}}
\toprule
\textbf{Provision} & \textbf{Variant} & \textbf{Change} & \textbf{Label} & \textbf{Annot.} & \textbf{XLM-R+fr+DA} & \textbf{XLM-R+DA} & \textbf{LaBSE-cal} \\
\midrule
rti-s11\_2 HI & pert. & number & 5.0 & 5 & -- & -- & -- \\
cpa-s10\_1 EN & pert. & negation & 2.5 & 1 & 3.2 & 4.8 & 7.6 \\
rti-s11\_3 HI & pert. & number & 5.0 & 5 & -- & -- & -- \\
cpa-s10\_2 EN & pert. & negation & 2.5 & 1 & 2.5 & 6.8 & 7.7 \\
rti-s11\_4 HI & pert. & modality & 5.0 & 5 & 8.5 & 8.3 & 7.7 \\
cpa-s12 EN & pert. & scope & 5.0 & 5 & 6.0 & 6.9 & 7.9 \\
rti-s12\_1 HI & pert. & modality & 5.0 & 5 & 10.0 & 8.3 & 7.9 \\
cpa-s14\_1 EN & pert. & negation & 2.5 & 1 & 2.6 & 7.1 & 7.5 \\
rti-s12\_4 HI & pert. & scope & 5.0 & 5 & 10.0 & 7.9 & 7.8 \\
cpa-s33\_1 EN & pert. & modality & 5.0 & 5 & 8.8 & 9.1 & 7.7 \\
rti-s14\_2 HI & pert. & exception drop & 7.0 & 9 & 8.0 & 9.5 & 7.7 \\
rti-s16\_1 HI & pert. & exception drop & 5.0 & 5 & -- & -- & -- \\
rti-s11\_2 HI & aggr. & omission & 6.5 & 4 & -- & -- & -- \\
rti-s11\_3 HI & aggr. & condition/ omission & 5.5 & 3 & -- & -- & -- \\
rti-s12\_6 HI & norm. & terminology & 9.5 & 9 & -- & -- & -- \\
rti-s14\_1 HI & norm. & obligation/ permission & 9.5 & 9 & 10.0 & 7.2 & 7.7 \\
rti-s15\_5 HI & aggr. & omission/ scope & 6.5 & 4 & 7.4 & 7.8 & 7.7 \\
cpa-s51\_2 EN & aggr. & exception/ omission & 8.5 & 8 & 7.9 & 9.6 & 7.2 \\
\bottomrule
\end{tabularx}
\end{table*}

{\footnotesize
\noindent\textbf{rti-s11\_2 (HI, perturbed).} Original: “\dev{जहां उपधारा (1) के अधीन, यथास्थिति, केन्द्रीय लोक सूचना अधिकारी या राज्य लोक सूचना अधिकारी द्वारा पर व्यक्ति पर किसी सूचना या अभिलेख या उसके किसी भाग के बारे में किसी सूचना की ताम}…”  Simplification: “\dev{जब केन्द्रीय या राज्य लोक सूचना अधिकारी उपधारा (1) के तहत किसी सूचना या अभिलेख के बारे में पर व्यक्ति को नोटिस भेजता है, तो उस पर व्यक्ति को नोटिस मिलने की तारीख से पंद्रह दिन के}…”  The third party has ten days from receipt of the notice, not fifteen.

\noindent\textbf{cpa-s10\_1 (EN, perturbed).} Original: “The Central Government shall, by notification, establish with effect from such date as it may specify in that notification, a Central Consumer Protection Authority to be known as…”  Simplification: “The Central Government will not set up, by notification and from a date it chooses, a Central Consumer Protection Authority, called the Central Authority. Its job is to deal with…”  Says the Central Government will NOT establish the Central Authority, reversing the provision and contradicting the rest of the sentence.

\noindent\textbf{rti-s11\_3 (HI, perturbed).} Original: “\dev{धारा 7 में किसी बात के होते हुए भी, यथास्थिति, केन्द्रीय लोक सूचना अधिकारी या राज्य लोक सूचना अधिकारी धारा 6 के अधीन अनुरोध प्राप्त होने के पश्चात् चालीस दिन के भीतर, यदि पर व्यक्}…”  Simplification: “\dev{धारा 7 में कुछ भी लिखा हो, अगर पर व्यक्ति को उपधारा (2) के तहत अपनी बात रखने का मौका दिया गया है, तो केन्द्रीय या राज्य लोक सूचना अधिकारी धारा 6 के तहत अनुरोध मिलने के तीस दिन के}…”  The forty-day period for the decision becomes thirty days.

\noindent\textbf{cpa-s10\_2 (EN, perturbed).} Original: “The Central Authority shall consist of a Chief Commissioner and such number of other Commissioners as may be prescribed, to be appointed by the Central Government to exercise the…”  Simplification: “The Central Authority will not be made up of a Chief Commissioner and a prescribed number of other Commissioners, all appointed by the Central Government to use the powers and car…”  Says the Central Authority will NOT consist of a Chief Commissioner and the prescribed Commissioners, reversing the composition provision.

\noindent\textbf{rti-s11\_4 (HI, perturbed).} Original: “\dev{उपधारा (3) के अधीन दी गई सूचना में यह कथन भी सम्मिलित होगा कि वह पर व्यक्ति, जिसे सूचना दी गई है, धारा 19 के अधीन उक्त विनिश्चय के विरुद्ध अपील करने का हकदार है।}”  Simplification: “\dev{उपधारा (3) के तहत भेजी गई सूचना में यह भी लिखा होगा कि जिस पर व्यक्ति को सूचना दी गई है, वह धारा 19 के तहत उस फ़ैसले के ख़िलाफ़ अपील करेगा।}”  The third party's entitlement to appeal (\dev{हकदार है}) becomes a statement that he will appeal (\dev{अपील करेगा}), turning a right into a duty.

\noindent\textbf{cpa-s12 (EN, perturbed).} Original: “No act or proceeding of the Central Authority shall be invalid merely by reason of— (a) any vacancy in, or any defect in the constitution of, the Central Authority; or (b) any def…”  Simplification: “No action or proceeding of the Central Authority will become invalid just because of: (a) a vacancy in, or a defect in how the Central Authority was formed; (b) a defect in the ap…”  Clause (c) covers 'any irregularity' not affecting the merits, not 'certain irregularity'.

\noindent\textbf{rti-s12\_1 (HI, perturbed).} Original: “\dev{केन्द्रीय सरकार, राजपत्र में अधिसूचना द्वारा, केन्द्रीय सूचना आयोग के नाम से ज्ञात एक निकाय का गठन करेगी, जो ऐसी शक्तियों का प्रयोग और ऐसे कृत्यों का पालन करेगा, जो उसे इस अधिनियम}…”  Simplification: “\dev{केन्द्रीय सरकार राजपत्र में अधिसूचना जारी करके केन्द्रीय सूचना आयोग नाम का एक निकाय बनाएगी। यह आयोग इस अधिनियम के तहत दी गई शक्तियों का इस्तेमाल करेगा और सौंपे गए काम कर सकेगा।}”  The Commission's duty to perform its assigned functions becomes optional ('\dev{सौंपे गए काम कर सकेगा}'), although the powers limb stays mandatory.

\noindent\textbf{cpa-s14\_1 (EN, perturbed).} Original: “The Central Authority shall regulate the procedure for transaction of its business and allocation of its business amongst the Chief Commissioner and Commissioners as may be specif…”  Simplification: “The Central Authority will not set its own procedure for carrying out its work and for dividing the work between the Chief Commissioner and the Commissioners, as laid down in regu…”  Says the Central Authority will NOT regulate its procedure and allocation of business, reversing the provision.

\noindent\textbf{rti-s12\_4 (HI, perturbed).} Original: “\dev{केन्द्रीय सूचना आयोग के कार्यों का साधारण अधीक्षण, निदेशन और प्रबंधन, मुख्य सूचना आयुक्त में निहित होगा, जिसकी सहायता सूचना आयुक्तों द्वारा की जाएगी और वह ऐसी सभी शक्तियों का प्रय}…”  Simplification: “\dev{केन्द्रीय सूचना आयोग के कामकाज की देखरेख, निर्देशन और प्रबंधन मुख्य सूचना आयुक्त के पास होगा, और सूचना आयुक्त उसकी मदद करेंगे। मुख्य सूचना आयुक्त आयोग की सभी शक्तियों का इस्तेमाल}…”  '\dev{ऐसे सभी कार्य}' (all such acts and things) becomes '\dev{कुछ काम}' (some acts), narrowing what the Chief Information Commissioner may do.

\noindent\textbf{cpa-s33\_1 (EN, perturbed).} Original: “The State Government shall provide the District Commission with such officers and other employees as may be required to assist the District Commission in the discharge of its func…”  Simplification: “The State Government may give the District Commission the officers and staff it needs to help it do its work.”  'Shall provide' the needed officers and employees becomes 'may give', turning a duty into a discretion.

\noindent\textbf{rti-s14\_2 (HI, perturbed).} Original: “\dev{राष्ट्रपति, उस मुख्य सूचना आयुक्त या सूचना आयुक्त को, जिसके विरुद्ध उपधारा (1) के अधीन उच्चतम न्यायालय को निर्देश किया गया है, ऐसे निर्देश पर उच्चतम न्यायालय की रिपोर्ट प्राप्त हो}…”  Simplification: “\dev{जिस मुख्य सूचना आयुक्त या सूचना आयुक्त के ख़िलाफ़ उपधारा (1) के तहत उच्चतम न्यायालय को मामला भेजा गया है, राष्ट्रपति उसे पद से निलंबित कर सकते हैं। ज़रूरी लगे तो जाँच के दौरान उसे}…”  Drops how long the suspension lasts, namely until the President passes orders on receiving the Supreme Court's report.

\noindent\textbf{rti-s16\_1 (HI, perturbed).} Original: “\dev{राज्य मुख्य सूचना आयुक्त .'ऐसी अवधि के लिए जो केन्द्रीय सरकार द्वारा विहित की जाए पद धारण करेगा और पुनर्नियुक्ति के लिए पात्र नहीं होगा : परन्तु कोई राज्य मुख्य सूचना आयुक्त पैंसठ}…”  Simplification: “\dev{राज्य मुख्य सूचना आयुक्त केन्द्रीय सरकार द्वारा तय की गई अवधि तक पद पर रहेगा और उसे दोबारा नियुक्त नहीं किया जा सकेगा।}”  Drops the proviso barring the State Chief Information Commissioner from holding office after the age of sixty-five.

\noindent\textbf{rti-s11\_2 (HI, aggressive).} Original: “\dev{जहां उपधारा (1) के अधीन, यथास्थिति, केन्द्रीय लोक सूचना अधिकारी या राज्य लोक सूचना अधिकारी द्वारा पर व्यक्ति पर किसी सूचना या अभिलेख या उसके किसी भाग के बारे में किसी सूचना की ताम}…”  Simplification: “\dev{अगर किसी दूसरे व्यक्ति की जानकारी दी जानी है, तो उसे दस दिन में अपनी आपत्ति बताने का मौका मिलेगा।}”  Drops the information officer who serves the notice and the fact that the ten days run from receipt of that notice, so the time limit has no starting point.

\noindent\textbf{rti-s11\_3 (HI, aggressive).} Original: “\dev{धारा 7 में किसी बात के होते हुए भी, यथास्थिति, केन्द्रीय लोक सूचना अधिकारी या राज्य लोक सूचना अधिकारी धारा 6 के अधीन अनुरोध प्राप्त होने के पश्चात् चालीस दिन के भीतर, यदि पर व्यक्}…”  Simplification: “\dev{किसी दूसरे की जानकारी के मामले में अधिकारी को चालीस दिन में तय करना होता है कि जानकारी दी जाए या नहीं, और उस व्यक्ति को लिखकर बताना होता है।}”  Drops the override of section 7, the sub-section (2) opportunity condition and the receipt of the section 6 request as the start of the forty days.

\noindent\textbf{rti-s12\_6 (HI, normal).} Original: “\dev{मुख्य सूचना आयुक्त या कोई सूचना आयुक्त, यथास्थिति, संसद् का सदस्य या किसी राज्य या संघ राज्यक्षेत्र के विधान-मंडल का सदस्य नहीं होगा या कोई अन्य लाभ का पद धारित नहीं करेगा या किसी}…”  Simplification: “\dev{मुख्य सूचना आयुक्त या कोई सूचना आयुक्त संसद का या किसी राज्य या संघ राज्यक्षेत्र की विधानसभा का सदस्य नहीं होगा। वह कोई दूसरा लाभ का पद नहीं रखेगा, किसी राजनीतिक दल से नहीं जुड़ेग}…”  All disqualifications are kept, but '\dev{विधान-मंडल}' (Legislature) is narrowed to '\dev{विधानसभा}' (Legislative Assembly).

\noindent\textbf{rti-s14\_1 (HI, normal).} Original: “\dev{उपधारा (3) के उपबंधों के अधीन रहते हुए, मुख्य सूचना आयुक्त या किसी सूचना आयुक्त को राष्ट्रपति के आदेश द्वारा साबित कदाचार या असमर्थता के आधार पर उसके पद से तभी हटाया जाएगा, जब उच्}…”  Simplification: “\dev{उपधारा (3) के प्रावधानों के अधीन, मुख्य सूचना आयुक्त या किसी सूचना आयुक्त को साबित कदाचार या असमर्थता के आधार पर केवल राष्ट्रपति के आदेश से पद से हटाया जा सकेगा। ऐसा तभी होगा जब र}…”  '\dev{हटाया जाएगा}' becomes '\dev{हटाया जा सकेगा}', a modality shift, though the exclusive route via the President's order and the Supreme Court report is preserved.

\noindent\textbf{rti-s15\_5 (HI, aggressive).} Original: “\dev{राज्य मुख्य सूचना आयुक्त और राज्य सूचना आयुक्त विधि, विज्ञान और प्रौद्योगिकी, समाजसेवा, प्रबंध, पत्रकारिता, जनसंपर्क माध्यम या प्रशासन और शासन में व्यापक ज्ञान और अनुभव वाले समाज}…”  Simplification: “\dev{राज्य सूचना आयुक्त समाज के जाने-माने लोग होने चाहिए। उन्हें कानून, विज्ञान, पत्रकारिता या सरकार चलाने जैसे कामों की अच्छी समझ होनी चाहिए।}”  Drops the State Chief Information Commissioner and several qualifying fields, and turns the closed list of fields into an open 'such as' list.

\noindent\textbf{cpa-s51\_2 (EN, aggressive).} Original: “Save as otherwise expressly provided under this Act or by any other law for the time being in force, an appeal shall lie to the National Commission from any order passed in appeal…”  Simplification: “If a State Commission decides an appeal, you can go to the National Commission next. It will only take the case if there is an important legal question.”  Drops the saving clause for anything expressly provided under this or another law, and 'important legal question' replaces substantial question of law.

}
Some trends are apparent. For inversion of negations across the entire operational provision, there is a discrepancy in annotations. One annotator considered these as off-track, while another considered them imprecise. The mean label (around 2.5) understates how misleading the text is: telling the reader that the Central Government “will not” establish the authority that the section establishes is telling them the opposite of the law. Conversely, changes in quantifiers and single digits received good agreement between annotators. This is because they are easily noticeable once identified. In practice, this means that the cases a deployed metric most needs to detect reliably are exactly those where annotators (human or in this case models) are least consistent.

Predicted values are based on E1 random splits, averaged across seeds where this pair appeared in test data. – indicates this pair did not appear in any E1 test splits. They overturn an expectation. While 3 inversion of negations (label 2.5) are predicted well by the main model (scores 3.2, 2.5, 2.6), versus 4.8-7.1 from scratch and 7.5-7.7 from calibrated LaBSE, the model performs poorly on other examples. Two Hindi Right to Information examples (label 5.0) receive perfect scores (10.0) for modality inversion (rti-s12\_1) and scope change (rti-s12\_4), a perfect score for a simplification that changes who must act or to whom the provision applies. Likewise, the English modality inversion (cpa-s33\_1) receives a high score (8.8). This corresponds to the diagnostic table, as modality is the worst category for the main model. The over-scoring failure JUDGEBERT identified in generic metrics thus survives in our learned metric, concentrated on obligation and permission and on Hindi. Good average performance does not prevent this issue. Therefore, rather than any correlation we report, we consider the diagnostic set to be the contribution most likely to outlast this paper.

\section{Reproducibility Notes}

Two defects in our own pipeline produced plausible-looking but worthless results, and both are worth reporting because neither announces itself.

The first is a tokenizer fallback. With transformers 5.x, if CamemBERT-v2 is loaded using the generic auto-tokeniser, a ‘byte-level’ tokenizer is used (without merge table), and words are split into individual characters. Sentence pairs become longer than 650 tokens. As a result, they are cut short to max-length and training degenerates to just guessing the mean label. Training error = variance in labels, test correlation slightly negative. No errors raised. Our loader now detects this case and falls back to the checkpoint's own tokenizer file, and we verify the token count for any new backbone.

The second is the Hindi text layer described in Section III-C. Both extractors failed in different ways. Although both outputs seem to be Hindi, they are erroneous. In general, explicit quantitative criteria should be applied to extraction stages in legal NLP tasks in Indic scripts because downstream metrics are computed happily on corrupted input.

End-to-end test of training code against published result: we converted the released FrJUDGE splits, and the original no data augmentation setting was re-run with the same 10 seeds on a single T4. This test was successful. However, note that by following the released code (100 epochs, best dev checkpoint), Pearson = 0.707 (\ensuremath{\pm} 0.042) and RMSE = 1.788 (\ensuremath{\pm} 0.115). These values differ from the published values (Pearson = 0.74 \ensuremath{\pm} 0.02, RMSE = 1.72). Our Pearson result is 0.033 lower. Given there are 10 seeds in each case, this is approximately 2 standard errors. Therefore although this confirms that the same approach has been taken, it is not exact replication. The most likely cause of difference from this point is due to versions of libraries in use (the tokenizer issue discussed below is an effect of transformers 5.x). Differences may also stem from our max token length (512) and non-deterministic outputs per seed in different hardware.

The more useful finding is that the setting matters more than the gap. With early stopping (patience=5), as described in the text of the paper, Pearson = 0.606 (\ensuremath{\pm} 0.284). For one seed, early stopping resulted in very poor performance (-0.196 at epoch 5). For the other 9 seeds, the mean was 0.695. However, in the released code, early stopping is commented out. Therefore, if a user followed the method described in the text, they may be confused by the high variability in results, and may conclude that the method does not work. However, the variability in results actually belongs to the stopping rule. Both results are presented.

In a partially run data augmented model (2 seeds completed), we can get some insight into the main result. With data augmentation, but evaluation only on human annotated test pairs, we achieve Pearson = 0.869 (\ensuremath{\pm} 0.023), and RMSE = 1.276 (\ensuremath{\pm} 0.128). The paper achieved Pearson = 0.97 on the test set, which includes trivial identical and unrelated data augmentation pairs. Therefore, the model with data augmentation is genuinely much stronger than the model without data augmentation. Furthermore, the published result is partly inflated by the contents of the test set. This is why we do not include data augmentation and sanity pairs in our main results. The FrJUDGE data is released under a CC-BY-NC-SA 4.0 license. It has been used to perform this verification, and as pre-training data for the transfer experiment. However, we do not redistribute it.

\section{Limitations and Ethics}

\begin{itemize}
\item The LMP labels in this version are disclosed model judgements, not human legal judgements. Therefore, we cannot confirm that the model agrees with legal judgements (the main claim of the previous work). The planned 30-pair human validation set is the minimum needed to begin making such evaluations.
\item Agreement between LLMs (Pearson 0.723) is the correlation between two instances of the same model, and may reflect common pre-training rather than a well-defined construct. It is also not an upper bound on predicting their mean. One model in this pair also generated the simplifications, meaning no true blind evaluation. The 68 largest disagreements were not adjudicated.
\item By simply predicting the mean training label of the variant for test pairs, which uses answer-key information and is not a deployable system, the benchmark achieves Pearson values higher than any learned models (see the E1 table). Therefore, the overall correlation cannot be used to determine whether a model considers legal effects, or just recognises simplification styles/error patterns. The within-variant table is the more honest measure.
\item No manual checking was performed on the Hindi text against the source PDFs. Only spot checks were performed by an automatic agent to verify the results of OCR and rule-based corrections. There are 3 Hindi identical pairs in the sanity holdouts, but no Hindi unrelated pairs. Therefore, checks on unrelated pairs were only performed on English examples.
\item Only two acts are represented in this corpus. For Hindi, only one act is represented, because no official Hindi text of the Consumer Protection Act was available when the corpus was built (the Hindi text of the Bill as introduced has since been located, and is not used here). There are discrepancies between the two versions of Right to Information. The Hindi version required OCR processing. Both discrepancies and OCR have been noted, and affected provisions are dropped rather than edited.
\item The clause review that gates the corpus was done by a rule-based agent rather than a human. The simplifications were written by an agent rather than sampled from an API at a recorded temperature. Both are disclosed above. The first thing we would add is a human pass over the retained clauses and over a sample of simplifications.
\item OCR was applied to the Hindi text. While the text passed quality assurance tests, and the clause numbering was corrected, there may still be vocabulary errors due to OCR. The retained clauses have not yet been checked by a legal domain expert in Hindi.
\item E2 is restricted to the 35 bilingual Right to Information clauses. This means that about 126 training examples were used per split. This is insufficient to draw conclusions about cross-lingual transfer.
\item FrJUDGE pre-training was unstable. When re-run with the same seed value, the FrJUDGE score was 0.691 compared to 0.873, and all downstream scores were reduced. In light of this instability, results obtained by single run transfer should be read with that variance in mind, as our non-augmented result of 0.586 is subject to this variability. The MuRIL and error-type results are post-hoc follow-ups and are reported as such.
\item The error-type classifier was trained and evaluated on rule-injected errors. Data used for training was augmented according to the same rules. Its macro-F1 therefore measures recognition of those five types, and cannot be used to measure the error types made by real simplifiers (omissions and paraphrasing).
\item Rule-injected errors may be easier to detect than errors a model makes naturally, so results on perturbed data are presented separately.
\item Due to the size of the dataset, the development set was particularly limited. Stopping the training process, and choosing the best checkpoint, was based on about 35 annotated example pairs. To improve performance, hyperparameters (learning rate, batch size and stopping criteria) were altered based on early stopping failing on the development set, however this means that results are optimised to some extent for the development set. Therefore results should be read as optimistic by an unknown margin. Ideally, this problem would be solved by increasing the dataset size and including a separate tuning set. As it stands, however, we can accurately say that hyperparameters were intentionally altered once, and not repeatedly tuned in response to the test set.
\item The text of statutes comes from publicly available government sources. This dataset is a course research artefact, for assessing quality of simplification, and is not a validated legal instrument. Predictions of an LMP score do not constitute legal advice.
\end{itemize}

\section{Conclusion}

A bilingual dataset of simplifications of Indian statutory text was constructed to allow quantification of limitations in previous work. This includes three contrasting simplification styles, and a diagnostic subset with an injected legal error of a known type. Splits at the clause level were confirmed to be disjoint, and sanity checks were applied to provisions not referenced elsewhere. Training a learned metric solely on this data gives a very unstable result (Pearson r = 0.228, std. dev. = 0.288), and does not consistently outperform a generic sentence embedding baseline. However, pre-training on FrJUDGE and then fine-tuning on the same splits gives 0.576 \ensuremath{\pm} 0.102 (English 0.675, Hindi 0.308), and the results are more consistent across seeds. Transfer between acts is possible in one direction (0.388) but barely in the reverse (0.119). In the subset of 35 clauses, there is insufficient data to determine whether transfer is possible in either direction between languages.

In addition, we have two further analyses. The MuRIL model with the same FrJUDGE pre-training matches the main system overall, with a Pearson of 0.577 (\ensuremath{\pm} 0.068) across five seeds. It generalises better, reaching 0.466 in Hindi and 0.251 in the hard cross-act direction. By repeating pre-training using XLM-R, we see that transfer performance is dependent on the pre-training run. The type of legal mistake in a pair can be predicted. With injected-error augmentation, a MuRIL classifier of legal mistake types achieves macro-F1 of 0.789 (vs 0.509 for rules), and so this provides some interpretability, to the extent possible due to injected rule-based mistakes.

However, there are three caveats. The first is that because the data has been constructed, a lookup that uses the answer key to predict each variant's mean label outperforms all learned models. Only the within-variant correlations (0.31-0.58) indicate whether models pick up legal information. The second caveat is that the diagnostic dataset separates the learned metric from general metrics, as designed. While negation mistakes are handled well, mistakes where an obligation was changed to a permission were poorly predicted. None of the systems passed both sanity checks. Finally, the measurement apparatus is sound; in the released setting, the same code reproduces the published FrJUDGE number. However, in the scenario described in the text, results are unstable. In order to make these results useful, two changes are required. Firstly, our disclosed model annotations must be replaced by human legal judgements. Otherwise, regardless of any correlations found, the central claim of this line of research is not being tested. Secondly, the failure mode in section VI must be fixed rather than averaged away. In cases where a metric scores perfectly on reversed obligations, this is extremely concerning given the intended use of the metric. No overall correlations (including those presented in this work) would make this acceptable.

\begin{thebibliography}{99}
\bibitem{ref1} D. Beauchemin, M. Albert-Rochette, R. Khoury and P.-L. Deziel, “JUDGEBERT: Assessing Legal Meaning Preservation Between Sentences,” in Proc. EMNLP, 2025, pp. 92-118.
\bibitem{ref2} T. Kew, A. Chi, L. Vasquez-Rodriguez, S. Agrawal, D. Aumiller, F. Alva-Manchego and M. Shardlow, “BLESS: Benchmarking Large Language Models on Sentence Simplification,” in Proc. EMNLP, 2023.
\bibitem{ref3} T. Zhang, V. Kishore, F. Wu, K. Q. Weinberger and Y. Artzi, “BERTScore: Evaluating Text Generation with BERT,” in Proc. ICLR, 2020.
\bibitem{ref4} F. Feng, Y. Yang, D. Cer, N. Arivazhagan and W. Wang, “Language-agnostic BERT Sentence Embedding,” in Proc. ACL, 2022.
\bibitem{ref5} A. Conneau, K. Khandelwal, N. Goyal, V. Chaudhary, G. Wenzek, F. Guzman, E. Grave, M. Ott, L. Zettlemoyer and V. Stoyanov, “Unsupervised Cross-lingual Representation Learning at Scale,” in Proc. ACL, 2020.
\bibitem{ref6} S. Khanuja, D. Bansal, S. Mehtani, S. Khosla, A. Dey, B. Gopalan, D. K. Margam, P. Aggarwal, R. T. Nagipogu, S. Dave, S. Gupta, S. C. B. Gali, V. Subramanian and P. Talukdar, “MuRIL: Multilingual Representations for Indian Languages,” arXiv preprint arXiv:2103.10730, 2021.
\bibitem{ref7} S. Doddapaneni, R. Aralikatte, G. Ramesh, S. Goyal, M. M. Khapra, A. Kunchukuttan and P. Kumar, “Towards Leaving No Indic Language Behind: Building Monolingual Corpora, Benchmark and Models for Indic Languages,” in Proc. ACL, 2023.
\bibitem{ref8} K. Krippendorff, “Computing Krippendorff's Alpha-Reliability,” Departmental Papers (ASC), Annenberg School for Communication, University of Pennsylvania, 2011.
\bibitem{ref9} T. Scialom, P.-A. Dray, S. Lamprier, B. Piwowarski, J. Staiano, A. Wang and P. Gallinari, “QuestEval: Summarization Asks for Fact-based Evaluation,” in Proc. EMNLP, 2021.
\bibitem{ref10} M. Maddela, Y. Dou, D. Heineman and W. Xu, “LENS: A Learnable Evaluation Metric for Text Simplification,” in Proc. ACL, 2023.
\bibitem{ref11} D. Beauchemin, H. Saggion and R. Khoury, “MeaningBERT: Assessing Meaning Preservation Between Sentences,” Frontiers in Artificial Intelligence, vol. 6, 2023, doi:10.3389/frai.2023.1223924.
\bibitem{ref12} S. Paul, A. Mandal, P. Goyal and S. Ghosh, “Pre-trained Language Models for the Legal Domain: A Case Study on Indian Law,” in Proc. ICAIL, 2023, pp. 187-196.
\bibitem{ref13} A. Joshi, S. Paul, A. Sharma, P. Goyal, S. Ghosh and A. Modi, “IL-TUR: Benchmark for Indian Legal Text Understanding and Reasoning,” in Proc. ACL, 2024.
\end{thebibliography}

\end{document}
